\pdfoutput=1


\begin{abstract}
    %Large language models (LLMs) have demonstrated impressive capabilities in tasks requiring deep knowledge understanding and complex reasoning. 
    Recent research on Reasoning of Large Language Models (LLMs) has sought to further enhance their performance by integrating meta-thinking---enabling models to monitor, evaluate, and control their reasoning processes for more adaptive and effective problem-solving. 
    However, current single-agent work lacks a specialized design for acquiring meta-thinking, resulting in low efficacy. 
    To address this challenge, we introduce \textbf{Re}inforced \textbf{M}eta-thinking \textbf{A}gents (\method), a novel framework that leverages Multi-Agent Reinforcement Learning (MARL) to elicit meta-thinking behaviors, encouraging LLMs to \textit{think about thinking}. 
    \method~decouples the reasoning process into two hierarchical agents: a high-level meta-thinking agent responsible for generating strategic oversight and plans, and a low-level reasoning agent for detailed executions. 
    Through iterative reinforcement learning with aligned objectives, these agents explore and learn collaboration, leading to improved generalization and robustness. 
    Empirical results from single-turn experiments demonstrate that \method~outperforms single-agent RL baselines on complex reasoning tasks, including competitive-level mathematical benchmarks and LLM-as-a-Judge benchmarks. 
    Additionally, we further extend \method~to multi-turn interaction settings, leveraging turn-level ratio and parameter sharing to improve efficiency.
    Comprehensive ablation studies further illustrate the evolving dynamics of each distinct agent, providing valuable insights into how the meta-thinking reasoning process enhances the reasoning capabilities of LLMs.
    Our code can be found in \texttt{https://github.com/ziyuwan/ReMA-public}
\end{abstract}
